% 2. Une discussion sur les changements à apporter pour passer de la version prototype
% à la version commerciale, tel que mentionné dans le mandat. Incluant une discussion
% sur le microcontrolleur, les interfaces des capteurs et une discussion sur les réseaux
% sans �l.

\subsection{Choix du microcontrolleur}
Le ESP32 est fiable, facile d'utilisation est extrèmement facile à intégré sur un PCB. Cependant, le modèle actuel est vieux.
Il serait donc meilleur d'utilisé le ESP32-S3, lequel à besoin de moins de circuit externe pour fonctionner, coute 3\$ de l'unité
et marche tel quel sur un PCB, sans grand circuit externe.
Il supporte nativement tout les protocoles nécessaires, peut être facilement connecté à une antenne LoRa, à un mode de consommation
qui utilise quelques micro ampères et de plus, il supporte nativement le protocole USB, sans circuit externe.

D'autres microcontrolleurs moins chère avec des consommations encore plus faible existe, mais celui-ci viens avec l'avantage que les
acheteurs pourraient le reprogrammé pour ajouté leurs propre fonctionalitées aux stations métérologique.

Il peut être utilisé pour les deux microcontrolleurs sans problème.

\subsection{Choix d'interfaces capteurs}
Les interfaces des capteurs, comme dit précédament, seraient meilleurs s'ils utilisaient tous le même protocol.
Cependant, temps qu'ils fonctionne avec les protocoles supporté par le ESP32-S3, c'est bon. En effet, les distances
sont courtes, la vitesse d'échange de donnée n'a pas besoin d'être très élevé, il n'y a essentiellement aucun interférence s'il
est utilisé dans des environnement normal. Sinon, de simple cables avec un "grounding" vas réglé le problème.

Tout les capteurs choisis ont des protocoles supportées par le ESP32-S3. Ils sont donc utilisable.

\subsection{Choix de communications réseaux}
LoRa sera utilisé pour les communications entre les stations météorologiques et la station de base, avec une topologie en étoile où 
la station de base agit comme maître qui envoie des requètes et les stations météo sont en mode basse consommation lorsqu'elles ne transmettent pas ou ne 
reçoivent pas de requêtes. 

LoRa permet une communication sans fil sur plusieurs kilomètres grâce à la modulation 
\textit{chirp spread spectrum} (CSS), avec des débits typiques de quelques centaines de bit/s à plusieurs kbit/s selon le 
\textit{spreading factor} (SF7 à SF12). C'est largement suffisant pour des données météorologiques, transmises par trames courtes, peu
fréquament. 

L'un des principaux avantages de LoRa est sa très faible consommation énergétique, avec des courants de veille en microampère et des 
transmissions de courte durée permettant un retour rapide en mode \textit{sleep}. Ça permet aux stations de fonctionner de manière 
autonome sur batterie ou énergie solaire pendant de longues périodes. La station de base vas être capable de gérer un nombre 
élevé de stations météo, limité principalement par le débit et le temps d'occupation du canal.

De plus, contrairement à Bluetooth, qui coute de l'argent lors de vente de modules au commercial, ce n'est pas le cas pour LoRa.